% =============================================================================
\section{Related Work}
% =============================================================================

\subsection{Gulf of Envisioning in LLM Interaction}

Subramonyam et al.\ identify a \textit{Gulf of Envisioning} in LLM-based
interaction: users struggle to form mental models of what a model can or will
produce, making it hard to specify intent precisely enough for reliable
output~\cite{subramonyam2024gulf}.
This gap is especially severe for SLMs, which lack the instruction-following
robustness of frontier models~\cite{zamfirescu2023johnny}.
Zamfirescu-Pereira et al.\ show that non-expert users systematically fail to
craft effective prompts even for LLMs, underspecifying constraints, omitting
output-format expectations, and iterating without a principled
strategy~\cite{zamfirescu2023johnny}.
Easy\_SLM addresses the Gulf of Envisioning specifically in the SLM context by
externalizing goal structure, decomposing prompts into steps, and capturing
audience and format constraints as first-class UI elements.

\subsection{Prompt Engineering Tools and Scaffolding}

Several systems support prompt construction for LLMs.
PromptMaker~\cite{jiang2022promptmaker} provides a prototype-driven workflow
for NLP practitioners; PromptIDE and similar tools target developers.
PROPEL~\cite{strobelt2022interactive} offers an interactive prompt development
environment with output inspection.
For general users, tools like ChatGPT's ``custom instructions'' or
Anthropic's system-prompt editor expose partial scaffolding but require users
to already understand prompt structure.
Our work differs in two ways: we target \textit{novice} users of
\textit{local SLMs}, and we integrate scaffolding as a research
\textit{control}---not to help users write better prompts as an end goal,
but to hold prompting variance constant so that the SLM--LLM capability
comparison is interpretable.

\subsection{SLM Accessibility and Local Deployment}

Despite a growing ecosystem of quantized, edge-deployable models (Qwen,
Llama, Phi, Mistral), HCI research on SLM-specific user experience is sparse.
Studies of LLM-powered tools (e.g.~\cite{brachman2024enterprise}) characterize
knowledge-worker use patterns in cloud contexts; local deployment adds
constraints---model loading latency, context-length limits, output brittleness
---that change the user experience substantially.
Our Reddit study explicitly samples SLM-specific communities (\textit{r/LocalLLaMA},
\textit{r/SelfHosted}) and compares their task discourse with general LLM
communities (\textit{r/ChatGPT}, \textit{r/OpenAI}) to characterize where the
two use-practice spaces overlap.

\subsection{Comparative Evaluation of AI Models}

Prior work comparing models in user-facing settings typically varies only the
model while keeping the interface fixed~\cite{lee2022coauthor}, or uses
crowdsourced preference data (e.g., LMSYS Chatbot Arena~\cite{chiang2024chatbot}).
Neither approach controls for prompting-skill variance.
Easy\_SLM operationalizes a controlled within-subjects design in which the same
scaffolded input specification is submitted to both model classes, allowing
objective metric computation alongside user preference capture---a design we
believe is novel for the SLM--LLM comparison problem.

\subsection{Sustainable and Resilient Computing}

Crawford's CHI 2024 keynote calls for ``rematerializing AI''---attending to the
physical and energy costs of large model
inference~\cite{crawford2024rematerializing}.
Offline-first interfaces are increasingly relevant as infrastructure fragility
becomes visible~\cite{tully2025crowdstrike}.
Easy\_SLM contributes to this discourse by making the SLM/LLM trade-off
legible to users rather than leaving it implicit in deployment choices made
by technologists.
